\chapter{Emulation settings: options and state of the art}\label{chap:emu-settings}

Every setting in this chapter changes the sound. For each one we give the possible choices, the default (in bold in
the tables), what you will hear when you change it, and what reference players and emulators do. When our default
differs from theirs, we explain why.

We invented almost none of these values. Filter presets come from published measurements of real chips and real
boards, and we name their origin each time. The raw list of every setting, with its internal key, is in
chapter~\ref{chap:all-settings}.

\section{General settings}
\begin{reglages}
\regl{Default duration}{30~s, 1, 2, \textbf{3}, 5, 10, 15, 30~min}{Used only when neither the file nor a database
gives the duration. Game music players usually choose between two and three minutes; three minutes let at least one
full loop of most tracks play.}
\regl{Silence detection}{\textbf{none}, 3, 4, 5, 10, 15~s}{Game Music Emu and the foobar2000 plugins built on it
skip to the next track after a silence. We leave it off: many tracks have deliberate pauses, and a SID playing
samples can seem silent for a few seconds.}
\regl{Loop limit}{\textbf{none}, 1 to 5}{For formats that know their loop point (VGM, some modules). Without a limit,
the known or default duration applies.}
\regl{Fade-in, fade-out}{\textbf{none}, 1, 2, 3, 5, 10~s}{A fade-out of a few seconds softens the cut of looping
tracks; VGMPlay and most VGM players offer one. We do not force it, so as never to alter a track that has a real
ending.}
\regl{Gain}{50 to 200~\%, \textbf{100~\%}}{Applied after emulation. Above 100~\%, tracks that are already loud clip:
turn up the phone volume instead.}
\end{reglages}

\section{SID (Commodore 64)}
\capture{settings-sid}{SID settings in the live panel}

The SID is the chip where the state of the art has gone furthest, and where settings matter most: two 6581s from the
same factory do not have the same filter. The software reference is the reSID family: Dag Lem's reSID, used by VICE,
and reSIDfp, developed later in libsidplayfp (sidplayfp on the command line) and ported to Java by JSIDPlay2.

\begin{reglages}
\regl{Engine}{reSID, \textbf{reSIDfp}}{sidplayfp uses reSIDfp by default, VICE offers reSID. reSIDfp models the 6581
filter from its transistors and reproduces combined waveforms better; reSID is lighter.}
\regl{Default model}{\textbf{6581}, 8580}{Used when the file does not say. The vast majority of HVSC music was written
on the 6581; it is also libsidplayfp's default.}
\regl{Force model}{\textbf{no}, yes}{To hear a track on the other chip. Useful for comparing, rarely faithful.}
\regl{C64 model}{\textbf{PAL}, NTSC, old NTSC, Drean (PAL-N), PAL-M}{The clock changes pitch and tempo. Most European
music is PAL; recent files state their model.}
\regl{Force C64 model}{\textbf{no}, yes}{As for the chip model.}
\regl{Interpolation}{\textbf{none}, performance, quality}{Maps to reSID's sampling methods: fast sampling,
interpolation, full resampling. Full resampling, which VICE also offers, removes treble aliasing but costs much more.
On a phone the difference is rarely audible.}
\regl{Filter emulation}{no, \textbf{yes}}{Without the filter, filtered voices sound raw: wrong for almost all C64
music. The setting is mostly there to understand what the filter brings.}
\regl{6581 preset}{\textbf{custom}, 31 presets}{See below.}
\regl{8580 preset}{\textbf{custom}, 5 presets}{See below.}
\regl{reSID: 6581 bias}{$-1000$ to $+1000$~mV, \textbf{0}}{Shifts the whole cutoff curve. VICE sets this bias to
$+500$~mV by default (\texttt{SidResidFilterBias}); we keep reSID's original value and offer VICE's as a preset.}
\regl{reSIDfp: 6581 curve}{0 (light) to 100 (dark), \textbf{50}}{Same scale and default as sidplayfp's
\texttt{FilterCurve6581} (0.5).}
\regl{reSIDfp: 6581 range}{0 (dark) to 100 (light), \textbf{50}}{Same default as sidplayfp.}
\regl{reSIDfp: 8580 curve}{0 (light) to 100 (dark), \textbf{50}}{Same default as sidplayfp.}
\regl{Combined waveforms}{\textbf{average}, weak, strong}{The strength of waveforms obtained by enabling several at
once varies from chip to chip. sidplayfp offers the same three levels (\texttt{CombinedWaveforms}), with the same
default.}
\regl{6581 pulse + saw}{\textbf{normal}, strong}{Some chips give a much louder pulse + sawtooth combination. Try it
when a sound seems missing in a track you know.}
\regl{8580 digi boost}{\textbf{no}, yes}{The 8580 makes samples played through the volume almost inaudible.
Musicians added a resistor to the board; libsidplayfp offers the same correction (\texttt{digiBoost}).}
\regl{2SID / 3SID output}{\textbf{stereo}, mono}{Stereo: one chip left, one right, the third centre, as most players
do. Mono: all chips centred, as on a real machine.}
\regl{Samples lane (\$D418)}{\textbf{no}, yes}{Specific to \soft. Separates samples played through the volume into
their own lane; the three voices then show only their oscillators. Applied to the next track.}
\end{reglages}

\subsection*{Filter presets}
The 31 6581 presets and the 5 8580 presets come from JSIDPlay2's default configuration. They fall into three
families.
\begin{description}
  \item[Lightest to Darkest] seven reSID curves, from lightest to darkest. JSIDPlay2 defines them by a curve
        position; we convert it to a reSID bias.
  \item[VICE default] VICE's $+500$~mV bias, to get that emulator's sound.
  \item[Measured chips] 23 reSIDfp curves named after the person who measured the chip and its date code: Alankila,
        Lord Nightmare, ZrX, Trurl, Nata, Grue, KruLLo, Enigma. ``6581R3 4885'' means an R3 revision made in week 48 of
        1985.
\end{description}
For the 8580: three reSID curves (Light, Average, Dark) and two chips measured by Trurl.

Our advice: leave the custom setting at its defaults, which are sidplayfp's. If a track sounds muffled, try Light or a
Trurl chip; if it whistles, Dark. There is no ``right'' 6581: the composer heard their own.

\section{Super Nintendo (SPC, SNSF)}
\capture{settings-spc}{SNES settings}
\begin{reglages}
\regl{Engine}{snes\_spc, \textbf{ares}}{ares carries on the work of bsnes and higan, the accuracy reference; Mesen is
the other emulator at that level. blargg's snes\_spc, used by Game Music Emu and many players, is lighter but does
not render the phase-inverted stereo some games use to widen the sound. Applied to the next track.}
\end{reglages}
We offer no interpolation other than Gaussian. Snes9x, bsnes and Mesen offer cubic or sinc interpolation, which makes
the sound brighter; pleasant, but not the console's sound, whose S-DSP applies exactly that Gaussian filter.

\section{Yamaha FM chips}
\capturedouble{settings-ym2612-psg}{settings-ym3812}{OPN2 settings}{OPL settings}
The core of an FM chip is digital: emulated well, it gives the same result as the chip. What changes from one
machine to another is what comes after, the DAC and the output circuit. That is why each family has a platform
setting, which sets the chip, low-pass filter, DC blocking and stereo from published measurements. If you change one
of those settings, the platform switches to ``custom''.

\subsection*{YM2612 and YM3438 (OPN2, Mega Drive)}
\begin{reglages}
\regl{Core}{\textbf{Nuked}, ymfm, Genesis Plus GX}{Nuked-OPN2 is written from the silicon and is the reference;
Genesis Plus GX also includes it as an option. ymfm is MAME's core. Genesis Plus GX's own core is the fastest.}
\regl{Platform}{\textbf{Mega Drive model 1 (VA0--VA6)}, model 1 VA7, model 2, Genesis 3, bare YM3438, bare YMF276,
custom}{Early Mega Drives have a separate YM2612; from model 1 VA7, the chip is built into the main ASIC as a YM3438.
Genesis Plus GX offers the same YM2612 / YM3438 choice.}
\regl{Chip}{\textbf{YM2612}, YM3438, YMF276}{Set by the platform.}
\regl{DAC ladder effect}{no, \textbf{yes} (model 1)}{Crossover distortion of the separate YM2612, which boosts quiet
sounds. Nuked-OPN2 reproduces it in YM2612 mode. The YM3438 lacks it.}
\regl{Low-pass filter}{0 to 24\,000~Hz, \textbf{3390~Hz}}{Model 1 filters its output with a first-order RC circuit.
BlastEm uses 3390~Hz; Genesis Plus GX's one-pole filter, set to 0.6, cuts around 3.7~kHz. Model 2 filters more: a
second order around 2840~Hz, measured by consolemods.org.}
\regl{Slope}{\textbf{6}, 12~dB/oct}{First order for model 1, second order for model 2.}
\regl{DC blocking}{0 to 200~Hz, \textbf{5~Hz}}{The output capacitors remove the DAC's DC voltage; this is essential
with the ladder effect, which offsets every voice even when silent. The real value is undocumented, so we put it
below the audible band.}
\regl{Stereo width}{0 to 100~\%, \textbf{100~\%}}{FM voices are only left, right or centre; reducing the width
helps on headphones with very separated tracks.}
\end{reglages}

\subsection*{YM2151 (OPM)}
\begin{reglages}
\regl{Core}{Nuked, \textbf{ymfm}, old MAME}{Nuked-OPM is the most exact but also the most expensive of all our cores:
we leave it as an option and use ymfm, MAME's current core.}
\regl{Platform}{\textbf{YM3012} (stereo: X68000, arcade), YM3014 (mono), custom}{The floating-point DAC is emulated
by the cores. Boards have no explicit filter on the FM part (noted in jotego's arcade work), so we add none.}
\regl{Filter, slope, DC, stereo}{\textbf{0~Hz}, \textbf{12~dB/oct}, \textbf{0~Hz}, \textbf{100~\%}}{No filtering by
default, for lack of measurements.}
\end{reglages}

\subsection*{YM2203, YM2608, YM2610 (OPN)}
\begin{reglages}
\regl{Core}{\textbf{ymfm}, old MAME}{There is no Nuked core for these chips. ymfm is MAME's core.}
\regl{Platform}{\textbf{PC-88 / PC-98}, custom}{The output stages of PC-88 and PC-98 boards have not been measured;
we do not filter.}
\regl{Clock (TurboFM)}{1.75, \textbf{3.5}, 3.58, 4~MHz}{For ZX Spectrum TurboFM music, the board runs at 3.5~MHz.}
\end{reglages}

\subsection*{YM3526, YM3812, YMF262 (OPL)}
\begin{reglages}
\regl{Core}{\textbf{Nuked}, AdLibEmu, MAME}{Nuked-OPL3 is the reference today, used in particular by modern DOSBox
versions. Ken Silverman's AdLibEmu and MAME's core are fast.}
\regl{Platform}{AdLib (mono), Sound Blaster 1.x/2.0, SB Pro 1/2, \textbf{SB16}, custom}{Values taken from DOSBox
Staging, which tuned them on recordings of real cards: SB 1.x and 2.0, first-order low-pass at 12~kHz; SB Pro,
first order at 8~kHz; SB16, no filter. The AdLib card has not been measured.}
\end{reglages}

\subsection*{YM2413 (OPLL)}
\begin{reglages}
\regl{Core}{\textbf{Nuked}, EMU2413, MAME}{Mitsutaka Okazaki's EMU2413 is the most widespread core in MSX and Master
System players and emulators; Nuked-OPLL is more exact.}
\end{reglages}

\section{AY-3-8910, YM2149 and the Z80 processor}
\capture{settings-ay}{AY settings}
\begin{reglages}
\regl{Clock}{ZX Spectrum (1.7734~MHz), \textbf{Pentagon (1.75~MHz)}, Taganrog (3.5~MHz), Atari ST (2~MHz), Amstrad CPC
(1~MHz)}{Most ZX scene music archived on ZXArt was written on the Pentagon, the most widespread clone in Russia, hence
our default. Files that state their clock take precedence.}
\regl{Type}{\textbf{AY-3-8910}, YM2149}{The YM2149 has a 32-step envelope instead of 16 and a slightly different
volume table. Choose it for the Atari ST.}
\regl{Interpolation}{none, performance, \textbf{quality}}{Smooths the conversion to 48~kHz. Without it the sound is
rawer, which some prefer on the Spectrum.}
\regl{Layout}{\textbf{ABC}, ACB, BAC, BCA, CAB, CBA, mono}{The original ZX Spectrum 128 is mono; stereo boards imposed
ABC or ACB depending on the region. Fuse and Unreal Speccy also offer ABC, ACB and mono.}
\regl{Mixer}{A: 100~\% left; B: 60~\% both sides; C: 100~\% right}{The three-channel mixer sets each voice. The centre
at 60~\% keeps a balanced level in mono.}
\regl{Duty cycle}{channels: \textbf{none} to all; 1 to 99~\%, \textbf{50~\%}}{A ZXTune option that matches no
hardware: it replaces the square with a pulse. Use it out of curiosity.}
\regl{Z80 clock}{\textbf{3.5~MHz}, 3.5469~MHz (128K), 4~MHz (CPC), 7~MHz}{For AY files that contain the player
code.}
\regl{Interrupt duration}{1 to 100~T, \textbf{24~T}}{How long the interrupt stays active each frame. Change it only if
a player misses interrupts.}
\end{reglages}

\section{SAA1099}
\begin{reglages}
\regl{Clock}{\textbf{SAM Coupé (8~MHz)}, Creative Music System (7.16~MHz)}{The two machines that use it.}
\regl{Interpolation}{none, \textbf{performance}, quality}{As for the AY.}
\end{reglages}

\section{Roland MT-32}
\begin{reglages}
\regl{ROM folder}{to choose}{Munt, also used by DOSBox and ScummVM, needs images of the module's control and sound
memories. We recognise them by their contents, whole or in halves.}
\regl{Preferred model}{\textbf{automatic}, MT-32, CM-32L / LAPC-I}{The CM-32L has extra sounds later games use; the
first games were written for the MT-32.}
\regl{Oscilloscope lanes}{\textbf{9 parts}, 32 partials}{Parts match the instruments; partials, the LA synthesis
building blocks. The second view will mostly interest those studying the module.}
\end{reglages}

\section{What we do not offer, and why}
\begin{itemize}
  \item No equaliser or ``surround'': they transform the machine's sound.
  \item No non-Gaussian interpolation for the Super Nintendo, for the same reason.
  \item No invented filter for boards that have not been measured (OPM, OPN, AdLib): we would rather add nothing than
        add a guess. As soon as a measurement is published, we will include it.
\end{itemize}
